\chapter{Relational Data: Postgres, JPA, Flyway}
\label{ch:postgres}

Two of the services own a PostgreSQL database each: auth-service owns
\code{ts\_auth}, course-service owns \code{ts\_course}. This chapter
covers the three layers between a \code{@Service} method and bytes on
disk --- JPA/Hibernate, the connection pool, and Flyway --- and the
architectural rule that there are \emph{two} databases at all.

\section{Database-per-service}

\begin{decision}[two Postgres servers, not two schemas in one]
Each service owns its data exclusively; the only way to another service's
data is through its API or its events. This forbids the classic failure
mode of shared databases --- service A quietly joining service B's tables,
coupling their schemas forever. The project takes it further than
required: not just separate databases but separate server processes, so
even accidental cross-connection is impossible and each can be tuned,
backed up, or broken independently. Cost on this box: about 256\,MiB per
extra server. In a cloud account the same isolation is usually two
databases on one managed RDS instance --- isolation by credentials rather
than by process --- because managed instances are billed per server.
\end{decision}

\section{JPA and Hibernate}

JPA is the specification (annotations like \code{@Entity},
\code{@OneToMany}, the \code{EntityManager}); Hibernate is the
implementation Spring Boot auto-configures. Spring Data JPA sits on top,
deriving queries from repository method names
(\code{findByCourseIdAndStudentId}) and generating the boilerplate.

The configuration from \code{configurations/course-service.yml}:

\begin{lstlisting}[language=yaml]
spring:
  datasource:
    url: jdbc:postgresql://localhost:5432/ts_course   (*@\co{1}@*)
    username: postgres
    password: root                                    (*@\co{2}@*)
  jpa:
    hibernate:
      ddl-auto: none                                  (*@\co{3}@*)
    show-sql: true                                    (*@\co{4}@*)
  flyway:
    enabled: true                                     (*@\co{5}@*)
\end{lstlisting}

\begin{description}
  \item[\co{1}] The dev default. Compose overrides host to
  \code{postgres-course}; Kubernetes overrides identically via
  \code{SPRING\_DATASOURCE\_URL}. Note the k8s Service was deliberately
  \emph{named} \code{postgres-course} so the override is the same string
  in both environments.
  \item[\co{2}] Plaintext credentials are acceptable only because this
  file is the \emph{dev default}; in the cluster,
  \code{SPRING\_DATASOURCE\_USERNAME}/\code{PASSWORD} arrive from a
  Kubernetes Secret (Chapter~16) and outrank these.
  \item[\co{3}] The most important line in the block ---
  Section~\ref{sec:ddlauto}.
  \item[\co{4}] Every SQL statement is logged --- a teaching tool and an
  N+1-query detector in dev; turn it off where log volume costs money.
  \item[\co{5}] Schema management belongs to Flyway, not Hibernate.
\end{description}

\section{Why \code{ddl-auto: none}}
\label{sec:ddlauto}

Hibernate can generate schema from entities (\code{ddl-auto: update}),
and every tutorial uses it. Production systems do not, because
\code{update} is a one-way ratchet with no history: it adds columns but
never drops or renames safely; two developers' databases diverge
silently; there is no record of what changed when, and no way to review a
schema change before it hits production.

\section{Flyway: migrations as code}

Flyway treats schema as an ordered series of immutable SQL scripts,
applied exactly once each, tracked in a table
(\code{flyway\_schema\_history}) inside the target database:

\begin{lstlisting}
auth-service/src/main/resources/db/migration/
  V1__create_users.sql
  V2__create_refresh_tokens.sql
  V3__add_admin_user_invitations.sql
  V4__seed_default_admin.sql
\end{lstlisting}

At startup, Flyway compares the history table with the classpath scripts
and applies the missing ones, in version order, before Hibernate ever
initializes. The rules that make it work: never edit an applied
migration (each row stores a checksum --- edits are detected and refused);
only add new versions; and write migrations to be compatible with the
\emph{previous} application version when you deploy rolling
(Chapter~15): add-then-migrate-then-remove, never rename in one step.

\begin{pitfall}[works locally, CrashLoopBackOff in the cluster]
Symptom: a service boots on your machine but its pod restarts endlessly;
logs show a Flyway checksum mismatch or a failed migration. Cause: the
local database had been migrated by an older (or edited) script --- your
laptop's history table and the cluster's disagree. This is also why the
Kubernetes startup probe (Chapter~17) targets an endpoint that only
answers after the Spring context is up: ``migrations applied'' is part of
what \emph{ready} means. Fix by adding a corrective migration, never by
editing history.
\end{pitfall}

\section{The connection pool}

Opening a Postgres connection costs a TCP handshake, authentication, and
a server process fork --- far too much per request. HikariCP (Spring
Boot's default pool) keeps a small set of open connections that requests
borrow and return. Default maximum: 10.

\begin{hardware}
Pool sizing folklore says ``more connections, more throughput''; the
reality on a 2-core server is the opposite --- Postgres does work
\emph{per connection}, and the practical formula
(\(\approx\) cores $\times$ 2) suggests the default 10 is already
generous for two services sharing two cores. If Postgres ever shows
\code{too many clients}, lower the pools before raising
\code{max\_connections}.
\end{hardware}

\section{Outside the project}

A managed-cloud equivalent of this chapter's setup --- AWS RDS Postgres
with the same Flyway flow --- changes only the URL and the trust model:

\begin{lstlisting}[language=yaml]
spring:
  datasource:
    url: jdbc:postgresql://app-db.xxxxx.eu-central-1.rds.amazonaws.com:5432/ts_course
    username: ${DB_USER}        # from AWS Secrets Manager
    password: ${DB_PASSWORD}
\end{lstlisting}

Backups, replication and failover become the provider's problem ---
which is why Chapter~19's warning about self-hosted state matters.
The second industrial variant worth knowing: running Flyway as a
separate CI step (\code{flyway migrate} against the target) instead of
at application startup, so schema changes deploy independently of code
and a slow migration cannot eat the startup probe's budget.

\begin{quiz}
\begin{enumerate}
  \item Give two concrete failure modes of \code{ddl-auto: update} in a
  team setting that \code{none} + Flyway prevents.
  \item Why must an applied migration never be edited, and what does
  Flyway use to detect it?
  \item You need to rename a column without downtime under rolling
  deployments. Sketch the multi-release sequence.
  \item Defend and attack the two-servers decision: one argument each
  way, and the setting in which each wins.
\end{enumerate}
\end{quiz}
